
\topic{Medicina Funcional de Precisión}

\slidetitle{Medicina Funcional de Precisión}
\begin{itemize}
    \item Es un cambio de paradigma hacia un modelo proactivo.
    \item Se utiliza la biología de sistemas para identificar la causa raíz de las enfermedades crónicas.
    \item Se logra estudiando las interacciones dinámicas del cuerpo humano.
\end{itemize}

\slidetitle{Personalización del tratamiento}
\begin{itemize}
    \item Este enfoque utiliza métricas individuales de salud para diseñar estrategias únicas.
    \item La intervención se adapta a la biomasa molecular y al contexto ambiental de cada paciente.
    \item No existe una prescripción única para todos en la medicina de precisión.
\end{itemize}

\slidetitle{Los órganos no están aislados}
\begin{itemize}
    \item Entendemos que el cuerpo es una red de sistemas interdependientes.
    \item A este mapa de conexiones moleculares lo llamamos el interactoma.
    \item Una falla en un sistema puede alterar la salud de órganos distantes.
\end{itemize}

\slidetitle{Identificación de la causa raíz}
\begin{itemize}
    \item Buscamos el origen del desequilibrio biológico en lugar de solo ocultar síntomas.
    \item Se utilizan heurísticas clínicas para navegar la complejidad del paciente crónico.
    \item El objetivo es restaurar la función fisiológica normal del organismo.
\end{itemize}

\slidetitle{Medicina de Redes}
\begin{itemize}
    \item Aplicamos la ciencia de redes y la física estadística a la medicina humana.
    \item Reconocemos que ninguna molécula o producto génico actúa solo en la célula.
    \item Las enfermedades complejas se ven como perturbaciones en redes moleculares integradas.
\end{itemize}

\slidetitle{Optimización del periodo de salud}
\begin{itemize}
    \item El propósito es que los años adicionales de vida sean funcionales y vibrantes.
    \item Empoderamos al paciente proporcionándole datos objetivos sobre su propia biología.
    \item Buscamos maximizar la resiliencia del huésped frente al entorno.
\end{itemize}

\slidetitle{El Interactoma Humano}
\begin{itemize}
    \item Es el mapa completo de interacciones físicas y funcionales en la célula.
    \item Dentro de este mapa identificamos "módulos de enfermedad" específicos.
    \item Estos módulos agrupan componentes que, al fallar, generan patologías.
\end{itemize}

\slidetitle{Genes y enfermedades crónicas}
\begin{itemize}
    \item Los genes que gobiernan enfermedades crónicas suelen estar débilmente vinculados en la red.
    \item No buscamos un solo "gen culpable", sino la interrupción de un módulo funcional.
    \item Este enfoque explica por qué las enfermedades crónicas son tan complejas.
\end{itemize}

\slidetitle{Fallos en cascada}
\begin{itemize}
    \item El cuerpo opera como una "red de redes" interdependientes.
    \item El fracaso de un nodo clave en una red puede provocar fallos en otros sistemas.
    \item Esto justifica la presencia de múltiples comorbilidades en el paciente inflamado.
\end{itemize}

\slidetitle{Concepto de Pathophenotype}
\begin{itemize}
    \item Clasificamos las enfermedades por sus mecanismos biológicos compartidos.
    \item Muchas patologías distintas pueden compartir la misma vía de inflamación sistémica.
    \item Tratamos el mecanismo de red subyacente y no solo la etiqueta del órgano.
\end{itemize}

\slidetitle{Predicción de respuestas celulares}
\begin{itemize}
    \item El modelado dinámico permite predecir cómo responderá un sistema a un cambio.
    \item Evaluamos la respuesta a fármacos o nutrientes según la topología de la red.
    \item Esto reduce los efectos secundarios y mejora la eficacia terapéutica.
\end{itemize}

\slidetitle{Individualidad bioquímica}
\begin{itemize}
    \item Dejamos atrás los promedios poblacionales para centrarnos en el paciente único.
    \item El individuo sirve como su propio control a través del seguimiento longitudinal.
    \item Adaptamos la medicina a la singularidad de la experiencia biológica humana.
\end{itemize}

\slidetitle{Metabolómica: la firma biológica}
\begin{itemize}
    \item Los metabolitos reflejan las actividades funcionales de la célula en tiempo real.
    \item Son el resultado final de la interacción entre nuestros genes y el ambiente.
    \item Nos dan una imagen dinámica y actualizada del estado de salud actual.
\end{itemize}

\slidetitle{Clasificación por Metabotipos}
\begin{itemize}
    \item Agrupamos a los individuos según sus perfiles metabólicos similares.
    \item Esto explica por qué dos pacientes con la misma enfermedad responden de forma distinta.
    \item Es la base para la nutrición y la farmacología personalizada de precisión.
\end{itemize}

\slidetitle{Detección con perfiles iPOP}
\begin{itemize}
    \item Utilizamos perfiles ómicos personales integrativos para monitorear al paciente.
    \item Capturamos cambios dinámicos en la sangre mucho antes de que aparezcan síntomas.
    \item Se han detectado riesgos de diabetes hasta una década antes de la falla glucémica.
\end{itemize}

\slidetitle{Señales en etapas subclínicas}
\begin{itemize}
    \item Interpretamos valores de laboratorio que se encuentran en los extremos de lo "normal".
    \item Estas fluctuaciones tempranas indican una pérdida de resiliencia homeostática.
    \item Intervenimos en etapas pre-patológicas para evitar el daño irreversible.
\end{itemize}

\slidetitle{Metabolitos y Microbioma}
\begin{itemize}
    \item Las bacterias intestinales producen mensajeros químicos como los ácidos grasos cortos.
    \item Estos metabolitos regulan la inflamación sistémica y la salud de la barrera hematoencefálica.
    \item Son piezas fundamentales para estabilizar el terreno biológico del huésped.
\end{itemize}

\slidetitle{IA y Ciencia de Datos}
\begin{itemize}
    \item La inteligencia artificial procesa la enorme complejidad de los datos ómicos.
    \item Identifica patrones no lineales en las redes metabólicas del paciente.
    \item Convierte el "Big Data" molecular en conocimiento clínico accionable.
\end{itemize}

\slidetitle{El Exposoma acumulado}
\begin{itemize}
    \item Es la suma de todas las exposiciones ambientales que recibimos desde la gestación.
    \item El 90\% del riesgo de enfermedades crónicas proviene del entorno y no de la genética.
    \item Incluye la dieta, tóxicos químicos, estrés psicológico y radiación.
\end{itemize}

\slidetitle{Carga tóxica y tejido adiposo}
\begin{itemize}
    \item Muchos xenobióticos del ambiente se almacenan en la grasa corporal.
    \item Desde allí, mantienen encendida una respuesta inflamatoria de bajo grado constante.
    \item Estos contaminantes pueden dañar directamente la función mitocondrial.
\end{itemize}

\slidetitle{Obesógenos ambientales}
\begin{itemize}
    \item Ciertos químicos industriales alteran los circuitos neuronales del hambre.
    \item Interfieren con la señalización de la insulina y el comportamiento alimentario.
    \item Explican la dificultad de muchos pacientes para regular su peso corporal.
\end{itemize}

\slidetitle{Detoxificación de precisión}
\begin{itemize}
    \item Evaluamos la capacidad genética del hígado para procesar y eliminar tóxicos.
    \item Cada paciente tiene una eficiencia distinta en sus vías de limpieza enzimática.
    \item Diseñamos planes de soporte nutricional basados en esta individualidad metabólica.
\end{itemize}

\slidetitle{Estabilización del Terreno}
\begin{itemize}
    \item Limpiar el terreno biológico es esencial para recuperar la capacidad de autorreparación.
    \item Mitigar el impacto del exposoma permite que las células vuelvan a su equilibrio.
    \item Un terreno resiliente es la base para el éxito de cualquier terapia avanzada.
\end{itemize}

\slidetitle{Epigenética: el archivo biológico}
\begin{itemize}
    \item El entorno y el estilo de vida modifican cómo se expresan nuestros genes.
    \item La metilación del ADN funciona como un historial de nuestras experiencias acumuladas.
    \item Estas marcas determinan qué genes se activan para protegernos o dañarnos.
\end{itemize}

\slidetitle{Relojes Epigenéticos}
\begin{itemize}
    \item Podemos medir la edad biológica real mediante biomarcadores de metilación.
    \item Esta edad a menudo no coincide con el tiempo cronológico del paciente.
    \item Un reloj acelerado indica un organismo bajo estrés biológico significativo.
\end{itemize}

\slidetitle{Reversibilidad epigenética}
\begin{itemize}
    \item Las marcas epigenéticas no son permanentes y pueden ser modificadas.
    \item Intervenciones en nutrición y manejo del estrés pueden frenar el envejecimiento celular.
    \item El paciente tiene el poder de reprogramar su propia trayectoria de salud.
\end{itemize}

\slidetitle{Programación en el útero}
\begin{itemize}
    \item El ambiente materno durante el embarazo deja huellas en el epigenoma del hijo.
    \item Un estado inflamatorio prenatal programa riesgos metabólicos para la vida adulta.
    \item Intervenir de forma temprana protege la salud de las futuras generaciones.
\end{itemize}

\slidetitle{Modulación con Ozonoterapia}
\begin{itemize}
    \item El ozono médico actúa como un fármaco que activa vías epigenéticas de defensa.
    \item Estimula al factor Nrf2 para potenciar la producción de antioxidantes endógenos.
    \item Reconfigura la respuesta celular frente a la inflamación crónica de forma duradera.
\end{itemize}

\slidetitle{Seguimiento Clínico Longitudinal}
\begin{itemize}
    \item Monitoreamos al paciente a través del tiempo mediante telemetría y sensores digitales.
    \item Esto nos permite ver la película completa de su salud y no solo fotos aisladas.
    \item Ajustamos el tratamiento de forma dinámica basándonos en datos en tiempo real.
\end{itemize}

\slidetitle{La Medicina del Futuro}
\begin{itemize}
    \item La medicina de precisión integra redes, ómicas y la salud del terreno biológico.
    \item Es una ciencia centrada en el ser humano para maximizar su potencial biológico.
    \item Estamos reescribiendo la historia de la salud para una longevidad funcional.
\end{itemize}

\sectiontitle{Listado de referencias}
\begin{itemize}
    \item Furman, D., et al. (2019). Nature Medicine, 25(12), 1822-1832.
    \item Loscalzo, J., et al. (2017). Network Medicine. Harvard University Press.
    \item Horvath, S., & Raj, K. (2018). Nature Reviews Genetics, 19(6), 371-395.
    \item Minich, D. M., & Bland, J. S. (2013). The Scientific World Journal.
    \item Cryan, J. F., et al. (2019). Physiological Reviews, 99(4), 1877-2013.
    \item García Vázquez, V. E. (2022). Ozonoterapia Médica. Cátedra de Medicina Traslacional.
\end{itemize}